% chapters/15/section-02.tex
\section{AudioClip, AudioSource, dan Ruang}
\label{sec:audioclip-audiosource}

\subsection{Dua Hal yang Berbeda}

Pada Unity, bunyi dipisahkan menjadi dua hal: berkasnya dan pemutarnya.
Perbedaan itu tampak sepele, tetapi menjelaskan sebagian besar kesulitan yang
akan Anda temui.

\begin{description}
  \item[\code{AudioClip}] Berkas bunyi yang telah diimpor, misalnya
        \code{lompat.wav}. Ia hanya berisi data. Ia tidak tahu di mana ia
        berada, dan ia tidak dapat berbunyi sendiri.
  \item[\code{AudioSource}] Komponen yang menempel pada sebuah
        \istilah{GameObject}. Ia yang memutar \code{AudioClip}, mengatur
        kekuatannya, dan menentukan dari mana bunyinya terdengar.
\end{description}

Membandingkannya dengan Bab~XIV berguna. \code{AudioClip} berperan seperti
\istilah{animation clip}: data yang tersimpan, dan tidak berbuat apa-apa
sendiri. \code{AudioSource} berperan seperti komponen \istilah{Animator}:
komponen yang menempel pada benda dan mengurus pemutarannya
\cite{unity-manual-import}. Karena itu,
pertanyaan ``di mana bunyi ini terdengar'' selalu dijawab oleh letak
\code{AudioSource}, bukan oleh berkasnya.

Rantai lengkapnya, dari peristiwa sampai ke pengeras suara, diperlihatkan pada
Gambar~\ref{fig:rantai-bunyi}. Kedua lajur bertemu pada
\istilah{AudioMixer}, dan itulah yang akan Anda pakai pada
Bagian~\ref{sec:waktu-volume}.

\begin{figure}[htbp]
\centering
% chapters/15/figures/rantai-bunyi.tikz
% Jalan bunyi dari peristiwa permainan sampai ke pendengar.
\begin{tikzpicture}[
  tahap/.style={
    rectangle, rounded corners=5pt, draw=ubPrimary, thick, fill=ubAlt,
    align=center, text width=24mm, inner sep=3.5pt, minimum height=11mm,
    font=\footnotesize
  },
  sumber/.style={
    rectangle, rounded corners=5pt, draw=ubAccent, thick, fill=white,
    align=center, text width=24mm, inner sep=3.5pt, minimum height=11mm,
    font=\footnotesize
  },
  panah/.style={-{Stealth[length=2.2mm]}, thick, draw=ubAccent},
  catatan/.style={font=\scriptsize\itshape, align=left},
  judul/.style={font=\footnotesize\bfseries, ubPrimary}
]

\node[judul, anchor=west] at (-1.35,3.55) {Lajur BGM (musik latar)};
\node[judul, anchor=west] at (-1.35,-0.05) {Lajur SFX (efek suara)};

% --- Lajur BGM ---
\node[tahap]  (mulai) at (0,2.5)    {scene dimuat};
\node[tahap]  (putar) at (3.2,2.5)  {\code{PutarMusik()}};
\node[sumber] (srcA)  at (6.4,2.5)  {\code{AudioSource}\\{\scriptsize\itshape loop = true}};

% --- Lajur SFX ---
\node[tahap]  (kejadian) at (0,-1.0)   {peristiwa\\{\scriptsize\itshape lompat, ambil, terluka}};
\node[tahap]  (pemicu)   at (3.2,-1.0) {\code{PemicuBunyi}};
\node[sumber] (srcB)     at (6.4,-1.0) {\code{AudioSource}\\{\scriptsize\itshape PlayOneShot}};

% --- Titik temu ---
\node[sumber] (mix)    at (9.6,0.75) {\code{AudioMixer}\\{\scriptsize\itshape Musik dan Efek}};
\node[tahap]  (keluar) at (12.8,0.75) {pengeras suara};

\draw[panah] (mulai) -- (putar);
\draw[panah] (putar) -- (srcA);
\draw[panah] (srcA) -- (mix);
\draw[panah] (kejadian) -- (pemicu);
\draw[panah] (pemicu) -- (srcB);
\draw[panah] (srcB) -- (mix);
\draw[panah] (mix) -- (keluar);

\node[catatan, anchor=north west, text width=64mm] at (-1.35,-2.0)
  {Kedua lajur bertemu pada \istilah{AudioMixer}. Di sanalah perbandingan
   kekuatan musik dengan efek suara diatur, tanpa perlu mengubah berkas
   bunyinya.};
\end{tikzpicture}

\caption{Jalan bunyi dari peristiwa permainan sampai ke pendengar}
\label{fig:rantai-bunyi}
\end{figure}

\subsection{Bunyi Datar dan Bunyi Beruang}

Setiap \code{AudioSource} memiliki pengaturan \istilah{Spatial Blend}, yaitu
penggeser dari nol sampai satu yang menentukan seberapa besar posisinya
berpengaruh.

\begin{itemize}
  \item \textbf{Nilai nol} berarti bunyi terdengar sama dari mana pun, seolah
        tidak memiliki tempat. Ini yang dipakai untuk musik latar dan seluruh
        bunyi antarmuka pengguna.
  \item \textbf{Nilai satu} berarti bunyi memiliki tempat di dalam ruang, dan
        kekuatannya berkurang menurut jarak pemain darinya. Ini yang dipakai
        untuk bunyi benda di dalam level.
\end{itemize}

Nilai di antaranya dipakai untuk keperluan khusus, misalnya bunyi yang perlu
sedikit memiliki tempat tetapi tidak boleh hilang ketika pemain menjauh.

Ketika \istilah{Spatial Blend} bernilai satu, jarak bunyi diatur melalui dua
nilai. \istilah{Min Distance} adalah jarak sampai bunyi mulai melemah; di
dalam jarak itu kekuatannya penuh. \istilah{Max Distance} adalah jarak saat
bunyi berhenti melemah lebih jauh. Nilai bawaan Unity sering kali terlalu
besar untuk permainan dua dimensi, sehingga bunyi-bunyi di sudut level masih
terdengar jelas padahal seharusnya sudah jauh.

\subsection{Play dan PlayOneShot}

Ada dua cara memanggil bunyi dari kode, dan pilihan di antara keduanya
menentukan perilaku yang Anda peroleh.

\begin{kodekeluaran}
// Cara pertama: satu sumber, satu bunyi pada satu waktu.
// Memanggilnya lagi memotong bunyi yang sedang berbunyi.
sfxSource.PlayOneShot(klipLompat);

// Cara kedua: mengganti isi sumber, lalu memutarnya.
// Bunyi sebelumnya berhenti, dan yang baru dimulai dari awal.
sfxSource.Stop();
sfxSource.clip = klipLompat;
sfxSource.Play();
\end{kodekeluaran}

Untuk efek suara, \code{PlayOneShot} hampir selalu pilihan yang tepat, sebab
ia mengizinkan beberapa bunyi berbunyi bersamaan pada satu
\code{AudioSource}. Memanggil \code{PlayOneShot} tiga kali dalam tiga bingkai
menghasilkan tiga bunyi yang bertumpuk. Sebaliknya, memakai \code{Play} dalam
keadaan yang sama menghasilkan satu bunyi yang selalu terpotong sebelum
selesai --- dan itulah yang terdengar seperti suara mesin yang patah-patah.

Untuk musik latar, \code{Play} yang dipakai, sebab hanya satu musik yang boleh
berbunyi pada satu waktu, dan kita justru menghendaki musik lama berhenti.

\begin{tipsbox}
\code{PlayOneShot} menerima satu argumen kedua yang jarang diketahui, yaitu
kekuatan bunyinya:

\begin{kodekeluaran}
sfxSource.PlayOneShot(klipLompat, 0.8f);
\end{kodekeluaran}

Argumen itu berguna ketika satu berkas bunyi dipakai untuk beberapa keperluan
dengan kekuatan yang berbeda, misalnya bunyi langkah yang lebih pelan ketika
pemain sedang menyelinap.
\end{tipsbox}
